Unmanned aerial vehicles (UAVs) are increasingly used in environmental monitoring, infrastructure inspection, emergency response, and other autonomous applications. As UAV systems become more complex, ensuring flight safety and reliability has become increasingly important. During operation, the flight-control system continuously records heterogeneous telemetry, including position estimates, inertial measurements, actuator responses, estimator innovations, battery states, and system-health variables. These signals provide valuable information about the runtime state of the vehicle and have therefore been widely used for UAV anomaly detection and fault diagnosis \cite{fourlas2021survey,puchalski2022uav,adaika2025fdd,tan2026anomaly,keipour2019automatic,hu2025mtcl,kabaoglu2026uavsead}.

Recent data-driven methods have shown strong capability in identifying abnormal flight states from multivariate telemetry. However, detecting an anomaly is only the first step toward practical diagnosis. Once abnormal behavior is observed, engineers still need to determine which fault domain is involved, which telemetry signals provide the strongest diagnostic evidence, and which software components should be inspected. Existing studies usually address these tasks separately: anomaly-detection methods focus on identifying abnormal states, explainable machine-learning methods provide feature-level attribution, and software fault-localization techniques operate on software-level evidence. The connection between telemetry-level diagnosis and flight-control software architecture therefore remains insufficiently explored \cite{tan2026anomaly,hu2025mtcl,lundberg2017shap,wong2016survey}.

UAV telemetry provides several properties that make such a connection feasible. Temporal-statistical features can summarize local variations in multivariate flight signals into compact and interpretable representations, while tree-based models are well suited to learning from these structured features \cite{wei2026aerotsboost,ke2017lightgbm}. Their predictions can be further explained using TreeSHAP, whose additive attributions can be aggregated from statistical features to physical telemetry channels and uORB topics \cite{lundberg2017shap,lundberg2020trees}. In PX4, uORB publisher--subscriber relationships explicitly connect telemetry topics with software modules, providing a natural architectural bridge from model evidence to software inspection \cite{meier2015px4,px4uorbdocs2026}.

To exploit these properties, we propose HS-AeroTS, a hierarchical framework for UAV telemetry anomaly diagnosis and architecture-aware software inspection. HS-AeroTS first performs runtime anomaly detection and conditional diagnosis over four fault domains: External Position, Global Position, Altitude, and Mechanical/Electrical. It then aggregates TreeSHAP contributions from temporal-statistical features to telemetry channels, uORB topics, and functional subsystems. Finally, the resulting topic-level evidence is mapped through a firmware-specific PX4 uORB dependency graph to generate software-module inspection candidates. This design connects anomaly detection, fault-domain diagnosis, interpretable telemetry evidence, and software architecture within a unified diagnostic pipeline.

We evaluate HS-AeroTS on the public UAV-SEAD benchmark using both chronological and flight-isolated settings \cite{kabaoglu2026uavsead}. The diagnostic evidence is further examined through attribution-based analyses and controlled PX4 replay experiments to assess its relevance and architecture-aware module ranking capability.

Our contributions are as follows:

\begin{itemize}[leftmargin=*,label=\textbullet]
    \item We propose HS-AeroTS, a hierarchical framework that extends UAV telemetry anomaly detection to fault-domain diagnosis and interpretable telemetry evidence.
    \item We develop a traceable evidence chain that aggregates TreeSHAP contributions from temporal-statistical features to telemetry channels and uORB topics, and further relates them to PX4 software-module inspection candidates through firmware-specific uORB relations.
    \item We evaluate the framework on the public UAV-SEAD benchmark and controlled PX4 replay, covering flight-level generalization, telemetry-evidence analysis, and architecture-aware software inspection.
\end{itemize}
